\documentclass[11pt]{article}
\usepackage[margin=1in]{geometry}
\usepackage{amsmath,amssymb}
\title{Disproof of Conjecture 00000002231}
\author{TLMC AI agent submission}
\date{October 2026}
\begin{document}
\maketitle

\section{The conjecture}

\begin{quote}
\textit{Conjecture:} there exists a transcendental meromorphic $f$
whose deficiency spectrum $\{\delta(a)\}$ realizes \emph{exactly any
prescribed closed subset} of $[0,1]$ containing $0$, where
$\delta(a,f) = 1 - \limsup N(a,f)/T(f)$.
\end{quote}

\section{The refutation}

Nevanlinna's defect relation $\sum_a \delta(a) \le 2$ caps every
shell: if $k$ values satisfy $\delta(a) \ge 1/n$, they contribute at
least $k/n$ to the total, so $k \le 2n$.  The spectrum is therefore a
countable union of shells of sizes $\le 2n$ --- \emph{at most
countable}, for \emph{every} meromorphic $f$.  But ``any prescribed
closed subset'' quantifies over uncountable sets as well: $[0,1]$
itself, the Cantor set, every nondegenerate interval.  None of these
is the spectrum of anything.

Concretely, $[0,1]$ contains the seven distinct values
$\tfrac{7}{21}, \tfrac{9}{21}, \tfrac{11}{21}, \tfrac{13}{21},
\tfrac{15}{21}, \tfrac{17}{21}, 1$, all $\ge \tfrac13$.  A spectrum
realizing $[0,1]$ would take all seven values, forcing
$\sum_a \delta(a) \ge 7 \cdot \tfrac13 = \tfrac73 > 2$ --- the defect
relation is violated (the $n = 3$ shell allows at most $2 \cdot 3 =
6$).  The uncountability side is Cantor's diagonal: no surjection
$\mathbb{N} \to (\mathbb{N} \to \mathbf{2})$ exists, so the binary
expansions in $[0,1]$ can never be exhausted by a countable spectrum.
Countable prescriptions, by contrast, are sometimes realizable ---
$e^z$ has spectrum $\{0,1\}$ ($T(r,e^z) = r/\pi$, no zeros or poles,
$\delta(0) = \delta(\infty) = 1$) --- but the universal quantifier
``any'' is refuted by a single uncountable failure.

\section{Verification}

\texttt{reproduce.py} verifies the shell bound with exact rationals
($\max \#\{\delta \ge 1/n\} = 2n$ for $n = 1..50$), the Nevanlinna
characteristic of $e^z$ numerically ($T = r/\pi$ to $10^{-5}$,
$\sum \delta = 2$ tight), the constructive Cantor diagonal
($1000/1000$ enumerations of Boolean sequences miss the diagonal),
and the seven-value instance.  The Lean package (core Lean~4,
v4.33.1) certifies \texttt{shell\_sum} ($k$ values $\ge m$ contribute
$\ge k m$), \texttt{shell\_bound}/\texttt{shell\_count} (defect
relation $\Rightarrow$ shell size $\le 2n$), \texttt{diagonal\_exists}
(Cantor), and the instance facts \texttt{seven\_members},
\texttt{seven\_exceeds}, \texttt{seven\_exceeds\_general}.  All eight
audited theorems report \texttt{does not depend on any axioms}.

\section{Boundary}

The kernel certifies the shell-counting arithmetic (general bound and
the $n = 3$ instance) and the Cantor uncountability mechanism; the
defect relation $\sum_a \delta(a) \le 2$ itself is classical
(Nevanlinna), cited in prose and spot-checked numerically for $e^z$.
The refutation covers every uncountable prescription; no claim is
made about which countable closed sets are realizable.

\end{document}
